% 2010 年湖北卷文科第 17 题：100 条鱼质量的频率分布直方图。
% 来源 PDF：普通高考/2010/2010湖北文.pdf，SHA-256 70d7537f9072187051a41095e97fab1d72cc357ee9cd04e831da8cf99217af78。
% 原页：第 1 页（A3 三栏）第 3 栏上部，书页码 599。
% 重绘依据：照原页——六个等宽组 [1.00,1.30)（组距 0.05），纵轴为「频率/组距」，
%   各矩形高依次为 0.4、1、4、5.6、6、3（矩形不加填充，故纵向虚线格线贯穿可见）；
%   纵轴刻度 0.4、1、3、4、5.6、6 处各引一条横向虚线指向对应矩形的顶边；
%   两轴均有断裂记号（横轴在 0 与 1.00 之间，纵轴贴近原点）。不代画任何数据。
\begin{tikzpicture}[x=1.5cm, y=0.72cm, semithick]
  % 坐标轴
  \draw[->, >=stealth] (0,0) -- (0,7.05);
  \draw[->, >=stealth] (-0.80,0) -- (6.95,0);
  % 轴断裂记号（横轴：0 与 1.00 之间；纵轴：贴近原点）
  \draw (-0.72,0.14) -- (-0.60,-0.14) -- (-0.48,0.14) -- (-0.36,-0.14);
  \draw (-0.09,0.08) -- (0.09,0.18);
  \draw (-0.09,0.24) -- (0.09,0.34);
  % 六个矩形（组距 0.05，高=频率/组距）
  \draw (0,0) rectangle (1,0.4);
  \draw (1,0) rectangle (2,1);
  \draw (2,0) rectangle (3,4);
  \draw (3,0) rectangle (4,5.6);
  \draw (4,0) rectangle (5,6);
  \draw (5,0) rectangle (6,3);
  % 纵轴刻度处的横向虚线：指向对应矩形的顶边
  % （0.4 一档即第一个矩形的顶边本身，不再另画）
  \draw[dashed, dash pattern=on 4pt off 3pt] (0,1) -- (1,1);
  \draw[dashed, dash pattern=on 4pt off 3pt] (0,3) -- (5,3);
  \draw[dashed, dash pattern=on 4pt off 3pt] (0,4) -- (2,4);
  \draw[dashed, dash pattern=on 4pt off 3pt] (0,5.6) -- (3,5.6);
  \draw[dashed, dash pattern=on 4pt off 3pt] (0,6) -- (4,6);
  % 纵轴刻度值
  \foreach \y/\v in {0.4/0.4, 1/1, 3/3, 4/4, 5.6/5.6, 6/6} {
    \node[anchor=east, inner sep=2pt] at (-0.16,\y) {$\v$};
  }
  % 横轴刻度值
  \node[anchor=north, inner sep=3pt] at (-0.72,-0.06) {$0$};
  \foreach \i/\v in {0/1.00, 1/1.05, 2/1.10, 3/1.15, 4/1.20, 5/1.25, 6/1.30} {
    \node[anchor=north, inner sep=3pt] at (\i,-0.06) {$\v$};
  }
  % 轴名
  \node[anchor=south, inner sep=1pt] at (-0.78,7.44) {频率};
  \draw (-1.30,7.14) -- (-0.26,7.14);
  \node[anchor=north, inner sep=1pt] at (-0.78,7.08) {组距};
  \node[anchor=west, inner sep=2pt] at (6.66,-0.34) {质量（$\mathrm{kg}$）};
\end{tikzpicture}
